% META: DONE.1

\section{Обсуждение ограничений и путей дальнейшего развития}

\subsection{Выявленные технические проблемы}
В ходе выполнения работы были идентифицированы две ключевые проблемы, ограничивающие точность и масштабируемость решения:

\begin{enumerate}
    \item \textbf{Качество и состав датасета}. Исходный датасет сильно несбалансирован по операциям: например, для \texttt{fillet} после всех фильтров осталось лишь 1~726 моделей, тогда как для \texttt{bossExtrusion} - значительно больше. Кроме того, многие модели используют сложные составные операции (например, массив отверстий), которые трудно представить как последовательность атомарных шагов. Это приводит к тому, что нейросеть обучается на упрощённых случаях и плохо обобщается на реальные промышленные детали.
    
    \item \textbf{Ограниченность STL как источника инженерной семантики}. STL описывает только геометрию поверхности, теряя информацию о гранях, эскизах, ограничениях и типах поверхностей. Это затрудняет различение, например, цилиндрического отверстия (созданного вырезанием) и сквозного паза (тоже вырезание, но другой формы). Переход к комбинированному использованию STL для распознавания и STEP/B-Rep для уточнения семантики является необходимым, но требует более сложной интеграции.
\end{enumerate}

\subsection{Перспективные направления развития}
На основе анализа результатов сформулированы следующие направления для улучшения системы (некоторые из них уже частично реализованы в прототипе):

\begin{enumerate}
    \item \textbf{Увеличение датасета за счёт использования промежуточных шагов}. Сейчас для обучения используется только последний шаг каждой модели, а если использовать все промежуточные состояния, размер датасета можно увеличить в 5-10 раз. Тем не менее, качество датасета может упасть ввиду добавления большого количество однотиповых моделей.
    
    \item \textbf{Переход к обучению на последних последовательных однотипных операциях}. Это позволит избежать штрафов за «сломанные» последующими операциями элементы (например, отверстие, перекрытое выдавливанием), а также позволит научиться естественным образом выделять порядок операций.
    
    \item \textbf{Улучшение постобработки}: добавление фильтрации по длине рёбер, по углу наклона, по расстоянию от центра модели для отсева выбросов.
    
    \item \textbf{Поддержка новых типов операций}: вращение (\texttt{bossRotated} и \texttt{cutRotated}), копия элементов по окружности (\texttt{circularCopy}), зеркальное отображение (\texttt{mirrorOperation}) и др. Для каждой новой операции потребуется собрать или сгенерировать соответствующий датасет.
    
    \item \textbf{Обучение на комбинированных датасетах}: сначала грубое обучение на большом «грязном» датасете, затем дообучение на малом, но тщательно сгенерированном и размеченном вручную.

    \item \textbf{Разделить обучение на прямолинейные и криволинейные контуры}: посколько обучение зависит от разбиения TODO

    \item \textbf{Улучшение триангуляции для нормализации числа треугольников}: TODO
    
    \item \textbf{Автоматическая оценка качества через КОМПАС-3D}: разработка скрипта, который измеряет процент успешно упрощённых моделей, суммарное число шагов восстановления и объём удалённой геометрии.
    
    \item \textbf{Увеличение размера нейронной сети и разрешения сетки} для улучшения распознавания мелких деталей.
\end{enumerate}

